\section*{अध्याय \ref{ch:g12:meiosis-genetic-shuffling} --- अर्धसूत्री विभाजन और आनुवंशिक फेंटाई}

\begin{solution}{exo:g12:meiosis-genetic-shuffling:1}
\omterm{def:g12:meiosis-genetic-shuffling:meiosis}{द्विगुणित}: हर \omterm{def:g11:cell-cycle-mitosis:chromosome}{गुणसूत्र} की दो प्रतियाँ ($2n$); \omterm{def:g12:meiosis-genetic-shuffling:meiosis}{अगुणित}: एक ($n$)। \omterm{def:g12:meiosis-genetic-shuffling:meiosis}{अर्धसूत्री
विभाजन} \omterm{def:g11:cell-cycle-mitosis:chromosome}{गुणसूत्र} संख्या आधी करता है, $2n$ से $n$ पर, और \omterm{def:g10:universal-dna:information}{DNA} को $2q$
(प्रतिकृति के बाद) से घटाकर प्रति युग्मक $q/2$ पर ले आता है।
\end{solution}

\begin{solution}{exo:g12:meiosis-genetic-shuffling:2}
अर्धसूत्री I हर जोड़े के समजात गुणसूत्र अलग करता है; और अर्धसूत्री II हर
\omterm{def:g11:cell-cycle-mitosis:chromosome}{गुणसूत्र} के दोनों \omterm{def:g11:cell-cycle-mitosis:chromosome}{क्रोमैटिड}। पहला न्यूनकारी है।
\end{solution}

\begin{solution}{exo:g12:meiosis-genetic-shuffling:3}
अर्धसूत्री I के बाद: 4 \omterm{def:g11:cell-cycle-mitosis:chromosome}{गुणसूत्र} (हर एक दो \omterm{def:g11:cell-cycle-mitosis:chromosome}{क्रोमैटिड} का), \omterm{def:g10:universal-dna:information}{DNA} $q$। अर्धसूत्री
II के बाद: 4 \omterm{def:g11:cell-cycle-mitosis:chromosome}{गुणसूत्र} (हर एक एक \omterm{def:g11:cell-cycle-mitosis:chromosome}{क्रोमैटिड} का), \omterm{def:g10:universal-dna:information}{DNA} $q/2$।
\end{solution}

\begin{solution}{exo:g12:meiosis-genetic-shuffling:4}
$2^4 = 16$।
\end{solution}

\begin{solution}{exo:g12:meiosis-genetic-shuffling:5}
ऐसे साथी से संकरण जो केवल \omterm{def:g11:genes-and-disease:genetic}{अप्रभावी} \omterm{def:g10:universal-dna:gene}{युग्मविकल्पी} ढोता हो, ताकि हर संतान
जाँचे गए जनक का युग्मक दिखा दे। पुनर्योजजों का अनुपात बता देता है कि दो \omterm{def:g10:universal-dna:gene}{जीन}
एक ही \omterm{def:g11:cell-cycle-mitosis:chromosome}{गुणसूत्र} पर हैं या नहीं और, यदि हैं, तो कितनी दूर।
\end{solution}

\begin{solution}{exo:g12:meiosis-genetic-shuffling:6}
अर्धसूत्री I: $2q$; विभाजनों के बीच: $q$; युग्मक: $q/2$। बीच में कोई S
प्रावस्था न होने से दूसरा विभाजन \omterm{def:g10:universal-dna:information}{DNA} फिर आधा कर देता है, और यही युग्मक को
किसी \omterm{def:g12:meiosis-genetic-shuffling:meiosis}{द्विगुणित} \omterm{def:g10:cells-common-unit:cell}{कोशिका} की आधी सामग्री पर ले आता है।
\end{solution}

\begin{solution}{exo:g12:meiosis-genetic-shuffling:7}
चार बराबर वर्ग: सहलग्न नहीं, अलग-अलग \omterm{def:g10:universal-dna:information}{गुणसूत्रों} पर। युग्मक $AB$, $Ab$,
$aB$, $ab$, हर एक चौथाई।
\end{solution}

\begin{solution}{exo:g12:meiosis-genetic-shuffling:8}
सहलग्न: दो बड़े वर्ग ($AB$, $ab$, पैतृक) और दो छोटे ($Ab$, $aB$,
पुनर्योजज)। पुनर्योजज $85/1000 = 8.5\%$।
\end{solution}

\begin{solution}{exo:g12:meiosis-genetic-shuffling:9}
केवल $AB$ और $ab$: दोनों \omterm{def:g10:universal-dna:gene}{जीनों} के ऊपर हुई कोई अदला-बदली ऐसे खंड बदलती है जो
इनमें से कोई नहीं ढोते, इसलिए $A/a$ और $B/b$ के संयोजन अनबदले रहते हैं। दो
\omterm{def:g10:universal-dna:gene}{जीनों} के बीच पुनर्योजन के लिए उनके बीच अदला-बदली चाहिए।
\end{solution}

\begin{solution}{exo:g12:meiosis-genetic-shuffling:10}
प्रति 1000 पर लगभग 0.8, 2.9 और 16: यानी 25 तथा 42 के बीच 20 का गुणक।
\end{solution}

\begin{solution}{exo:g12:meiosis-genetic-shuffling:11}
अर्धसूत्री I पर दोनों समजात एक ही \omterm{def:g10:cells-common-unit:cell}{कोशिका} में जाते हैं: दो युग्मक दो-दो
प्रतियों वाले और दो बिना किसी प्रति के। और अर्धसूत्री II पर किसी एक
\omterm{def:g11:cell-cycle-mitosis:chromosome}{गुणसूत्र} के दोनों \omterm{def:g11:cell-cycle-mitosis:chromosome}{क्रोमैटिड} एक ही \omterm{def:g10:cells-common-unit:cell}{कोशिका} में जाते हैं: एक युग्मक दो
प्रतियों वाला, एक बिना, और दो सामान्य।
\end{solution}

\begin{solution}{exo:g12:meiosis-genetic-shuffling:12}
पुनर्योजन दूरी के समानुपाती है: 2\% का मतलब वे \omterm{def:g10:universal-dna:gene}{जीन} बहुत पास हैं, इसलिए कोई
अदला-बदली शायद ही उनके बीच पड़े; और 48\% का मतलब बहुत दूर, जहाँ लगभग हर
\omterm{def:g12:meiosis-genetic-shuffling:meiosis}{अर्धसूत्री विभाजन} में उनके बीच अदला-बदली होती है। चूँकि एक अदला-बदली 50\%
पुनर्योजज युग्मक देती है, इसलिए दूर पड़े सहलग्न \omterm{def:g10:universal-dna:gene}{जीन} स्वतंत्र \omterm{def:g10:universal-dna:gene}{जीनों} जितनी ही
आज़ादी से पुनर्योजित होते हैं।
\end{solution}

\begin{solution}{exo:g12:meiosis-genetic-shuffling:13}
$AB$, $Ab$, $aB$, $ab$, हर एक चौथाई, दोनों जोड़ों के स्वतंत्र अपव्यूहन से।
\omterm{prop:g12:meiosis-genetic-shuffling:crossingover}{जीन-विनिमय} इन \omterm{def:g10:universal-dna:gene}{जीनों} के लिए कुछ नहीं बदलता: वह किसी एक \omterm{def:g11:cell-cycle-mitosis:chromosome}{गुणसूत्र} की लंबाई भर
\omterm{def:g10:universal-dna:gene}{युग्मविकल्पी} पुनर्योजित करता है, और ये अलग-अलग \omterm{def:g10:universal-dna:information}{गुणसूत्रों} पर हैं।
\end{solution}

\begin{solution}{exo:g12:meiosis-genetic-shuffling:14}
हर जनक हर संतान को हर \omterm{def:g10:universal-dna:gene}{जीन} पर दो \omterm{def:g10:universal-dna:gene}{युग्मविकल्पियों} में से एक देता है, जो
यादृच्छिक ढंग से चुना जाता है; और किसी दिए गए \omterm{def:g10:universal-dna:gene}{जीन} पर दो भाई-बहनों को किसी
जनक से वही \omterm{def:g10:universal-dna:gene}{युग्मविकल्पी} $1/2$ प्रायिकता से मिला। बहुत-से \omterm{def:g10:universal-dna:gene}{जीनों} पर साझा अंश
का औसत $1/2$ आता है, पर किसी एक भाई-बहन जोड़ी के लिए वह उसके इर्द-गिर्द
बदलता है।
\end{solution}

\begin{solution}{exo:g12:meiosis-genetic-shuffling:15}
मातृ पौधे के क्लोन, जो आनुवंशिक रूप से उसके और आपस में समरूप हैं। वह पौधा
ऐसे \omterm{def:g11:enzymes-and-phenotype:phenotype}{जीनप्ररूप} का तेज़ और पक्का गुणन पाता है जो यहाँ और अभी काम करता है; और
वह विविधता खोता है जो कुछ वंशजों को पर्यावरण के किसी बदलाव या किसी नए
परजीवी से बचने देती।
\end{solution}

\begin{solution}{pb:g12:meiosis-genetic-shuffling:1}
\textbf{1.} मादा $b^+b$, $vg^+vg$, जिसके एक \omterm{def:g11:cell-cycle-mitosis:chromosome}{गुणसूत्र} पर $b^+vg^+$ और दूसरे
पर $b\,vg$ है (उसके दोनों शुद्ध जनकों से)। नर $bb$,
$vg\,vg$।

\textbf{2.} केवल $b\,vg$ युग्मक: वह नर केवल \omterm{def:g11:genes-and-disease:genetic}{अप्रभावी} \omterm{def:g10:universal-dna:gene}{युग्मविकल्पी} देता है,
इसलिए हर संतान का \omterm{def:g11:enzymes-and-phenotype:phenotype}{लक्षणप्ररूप} उस मादा का युग्मक बता देता है।

\textbf{3.} $b^+vg^+$ (धूसर-सामान्य), $b\,vg$ (काली-अवशेषी),
$b^+vg$ (धूसर-अवशेषी), $b\,vg^+$ (काली-सामान्य)।

\textbf{4.} धूसर-सामान्य 42\%, काली-अवशेषी 41\%, धूसर-अवशेषी 9\%,
काली-सामान्य 8\%।

\textbf{5.} एक ही \omterm{def:g11:cell-cycle-mitosis:chromosome}{गुणसूत्र} पर: $1:1:1:1$ के बजाय दो बड़े और दो छोटे वर्ग।

\textbf{6.} धूसर-सामान्य और काली-अवशेषी: यानी वे संयोजन जो उसके दोनों समजात
ढोते थे, और जो उसके शुद्ध जनकों से अटूट विरासत में मिले।

\textbf{7.} धूसर-अवशेषी और काली-सामान्य, जो पूर्वावस्था I में दोनों \omterm{def:g10:universal-dna:gene}{जीनों} के
बीच हुए किसी \omterm{prop:g12:meiosis-genetic-shuffling:crossingover}{जीन-विनिमय} से आए।

\textbf{8.} $(206 + 185)/2300 \approx 17\%$। दोनों \omterm{def:g10:universal-dna:gene}{जीनों} के बीच का
\omterm{prop:g12:meiosis-genetic-shuffling:crossingover}{जीन-विनिमय} चारों में से दो \omterm{def:g11:cell-cycle-mitosis:chromosome}{क्रोमैटिड} पुनर्योजज बनाता है, इसलिए 17\%
पुनर्योजज युग्मकों का मतलब लगभग 34\% अर्धसूत्री विभाजनों में अदला-बदली।

\textbf{9.} एक अकेली अदला-बदली दोनों पुनर्योजज \omterm{def:g11:cell-cycle-mitosis:chromosome}{क्रोमैटिड} साथ-साथ बनाती है,
हर प्रकार का एक; और दोनों समजात युग्मकों में बराबरी से जाते हैं, इसलिए दोनों
पैतृक वर्ग भी मेल खाते हैं।

\textbf{10.} अब पैतृक वर्ग धूसर-अवशेषी और काली-सामान्य हैं, हर एक लगभग
41.5\%; और पुनर्योजज धूसर-सामान्य तथा काली-अवशेषी, हर एक लगभग
8.5\%।

\textbf{11.} हाँ: दोनों के साथ पुनर्योजन 50\% से काफ़ी नीचे।

\textbf{12.} $b$ --- $cn$ --- $vg$: 9 और 8 जुड़कर 17 बनते हैं।

\textbf{13.} किसी अंतराल में अदला-बदली की संभावना उसकी लंबाई के समानुपाती
है, इसलिए सटे हुए अंतरालों की बारंबारताएँ जुड़ जाती हैं: पुनर्योजन
बारंबारता \omterm{def:g11:cell-cycle-mitosis:chromosome}{गुणसूत्र} पर दूरी नापती है।

\textbf{14.} एक अदला-बदली चार में से दो \omterm{def:g11:cell-cycle-mitosis:chromosome}{क्रोमैटिड} पुनर्योजज देती है: यानी
अधिक से अधिक 50\%; और अतिरिक्त अदला-बदलियाँ फेंटती तो हैं पर पुनर्योजज अंश
कभी आधे से ऊपर नहीं उठातीं।

\textbf{15.} वह किसी और \omterm{def:g11:cell-cycle-mitosis:chromosome}{गुणसूत्र} पर है या उसी पर बहुत दूर; ये संकरण यह नहीं
बता सकते कि कौन-सा।

\textbf{16.} $2^4 = 16$ युग्मक; $16 \times 16 = 256$ संयोजन।

\textbf{17.} हर अदला-बदली ऐसे \omterm{def:g11:cell-cycle-mitosis:chromosome}{गुणसूत्र} बनाती है जो पहले कभी थे ही नहीं, और
वह भी ऐसी जगह पर जो एक \omterm{def:g12:meiosis-genetic-shuffling:meiosis}{अर्धसूत्री विभाजन} से दूसरे में बदलती है; इसलिए
अलग-अलग \omterm{def:g10:universal-dna:information}{गुणसूत्रों} की, यानी युग्मकों की, संख्या की कोई व्यावहारिक सीमा नहीं
है।

\textbf{18.} $2^{23} \times 2^{23} \approx \num{7e13}$: यानी अब तक जी चुके
मनुष्यों की संख्या से सात सौ गुना अधिक संयोजन, और वह भी \omterm{prop:g12:meiosis-genetic-shuffling:crossingover}{जीन-विनिमय} से पहले।

\textbf{19.} भाई-बहन अपने \omterm{def:g10:universal-dna:gene}{युग्मविकल्पी} उन्हीं चार समुच्चयों से निकालते हैं
(हर जनक से दो) और औसतन उनका आधा साझा करते हैं; जबकि अजनबी अलग-अलग समुच्चयों
से निकालते हैं।

\textbf{20.} 17\%, यानी उस मादा के युग्मकों का वह अंश जो $b$ और $vg$ के बीच
पुनर्योजित हुआ, दोनों \omterm{def:g10:universal-dna:gene}{जीनों} के बीच की दूरी नापता है; और जोड़ों का स्वतंत्र
अपव्यूहन तथा हर जोड़े के भीतर का \omterm{prop:g12:meiosis-genetic-shuffling:crossingover}{जीन-विनिमय} हर युग्मक को अलग बनाते हैं।
\end{solution}
